\pdfmapfile{+cm.map}
\pdfmapfile{+cmextra.map}
\pdfmapfile{+symbols.map}
\documentclass[11pt]{article}
\usepackage[margin=0.9in]{geometry}
\usepackage{amsmath,amssymb,amsthm,booktabs,longtable,array,hyperref}
\hypersetup{colorlinks=true,linkcolor=blue,urlcolor=blue,pdftitle={Fourteen integer unknowns and smaller norms for Con(ZFC)},pdfauthor={AI-assisted research for ProveIt}}
\newtheorem{theorem}{Theorem}
\newtheorem{lemma}{Lemma}
\newtheorem{proposition}{Proposition}
\newtheorem{corollary}{Corollary}
\newcommand{\Z}{\mathbb Z}
\newcommand{\N}{\mathbb N}
\newcommand{\PP}{\mathbb Z_{>0}}
\newcommand{\Con}{\operatorname{Con}}
\newcommand{\Sq}{\operatorname{Sq}}
\renewcommand{\labelitemi}{$\bullet$}
\setlength{\emergencystretch}{3em}
\title{Fourteen integer unknowns and smaller separating norms\\for the consistency of ZFC}
\author{AI-assisted research prepared for Vladimir Reshetnikov's ProveIt repository}
\date{23 October 2026}
\begin{document}
\maketitle
\begin{abstract}
We compare the seven recent Diophantine presentations of the blank computation of K215 and combine their strongest independent reductions. The principal new construction uses fourteen unrestricted integer unknowns. A Pell orbit removes one order constraint, and a cofinal choice of the Jones power-of-two base removes two more. We also obtain smaller twelve-positive-variable norms, a lower-degree nine-positive-variable equation, a fifty-integer quartic, and a cheaper moderate-degree twenty-four-integer equation. The smallest positive-variable form still has nine unknowns. Every equivalence below is an ordinary mathematical argument formalizable in ZFC, using the identified compiler and machine-calibration results. Finite replay checks supplement these arguments; they are not new end-to-end Lean or Rocq proofs. No consistency decision, global minimum, or universal representation record is claimed.
\end{abstract}

\section{Scope, baseline, and what ``simpler'' means}
The review is pinned to ProveIt commit
\texttt{8c6eec9d19ef80538c8e61ff6e0180040583446a}.
The source collection is
\href{https://github.com/VladimirReshetnikov/ProveIt/tree/8c6eec9d19ef80538c8e61ff6e0180040583446a/SetTheory/Cardinals/docs/reports/hilbert-tenth-problem/con-zfc-diophantine-forms}{Con(ZFC) in Diophantine form}, specifically packages 24--30.
The previous delivery is package 27. There is no single order on simplicity: domain, number of unknowns, equation count, total degree, arithmetic-circuit cost, coefficient height, and proof length must be distinguished. In particular, nine positive unknowns cannot simply be read as nine unrestricted integers.

All operation counts charge binary addition, subtraction, and multiplication, including multiplication by numerals and squaring. Reuse of a previously computed value and comparison of two system outputs are free. Computing the fixed coefficients is not charged to evaluation of a polynomial in its unknowns. Thus these counts are not bit-complexity estimates. Huge constants are supplied either as actual binary integers or exact bounded arithmetic recipes, never as unspecified existential parameters.

\begin{center}\small
\begin{tabular}{p{0.23\linewidth}p{0.24\linewidth}p{0.39\linewidth}}
\toprule
Source & Main contribution & Use in this synthesis\\\midrule
24, integer24 & Degree 1069 positive norm; sharp signed weights & Halved carrier and leading cancellation; signed weight bound\\
25, hybrid12 & Degree 1081 at 116 operations; smaller tuple & Comparison; whole-tuple discipline retained\\
26, ratio12 & Ratio-bounded roots and cheaper systems & Comparison of domains and free-comparison costs\\
27, previous delivery & Fast nine-variable circuits; degree 1081/120; integer lifts & Shared-power and coefficient-alias evaluation; signed and low-degree proofs\\
28, compact norms & Structured majorants and signed finalizer & Lower-degree nine-variable equation\\
29, simpler forms & Explicit 17-integer Jones specialization & Square stack, intersective merger, separated three-root norm\\
30, quartic31 & 31-variable quartic; trace/remainder norm; 111-operation product & Smaller Jones parameter; register elimination; scalar norm evaluator\\\bottomrule
\end{tabular}
\end{center}
The collection's corrected reconciliation is included as \texttt{sources/README.md}. It is an intake comparison, not an independent proof of every delivered assertion. Our comparison distinguishes identical polynomials from equal-looking circuits in different coefficient tuples. In particular packages 25 and 27 share two coefficient integers but not the entire tuple. Likewise replacing a tuple in package 18's circuit changes its fixed polynomial.

\subsection*{Selected new frontiers}
Put $L=5^{32}$. The selected equations below all satisfy, in their stated domains,
$\mathrm{ZFC}\vdash (\text{no solution})\leftrightarrow\Con(\mathrm{ZFC})$.
\begin{center}\small
\begin{tabular}{llrr}
\toprule
Domain and unknowns & Degree & Operations & Equations\\\midrule
$14\Z$ & $813056L+2752$ & 1524 & 1\\
$16\Z$ & $406528L+10008$ & 1537 & 1\\
$9\PP$ & $92324L+296$ & 1500 & 1\\
$9\N$ & $92324L+296$ & 1505 & 1\\
$12\PP$ & $1061$ & 120 & 1\\
$12\PP$ & $1065$ & 119 & 1\\
$12\PP$ & $1069$ & 117 & 1\\
$12\PP$ & $1081$ & 109 & 1\\
$12\PP$ & $1052$ & 114 & 3\\
$12\PP$ & $1072$ & 103 & 3\\
$24\Z$ & $1936$ & 161 & 1\\
$24\Z$ & $1920$ & 156 & 3\\
$50\Z$ & $4$ & 197 & 1\\
$50\Z$ & $2$ & 111 & 29\\\bottomrule
\end{tabular}
\end{center}
The two selected twelve-variable systems use an explicit transport-coordinate shift; their variable domain is unchanged. The 109-operation product is literally the source30 polynomial evaluated more cheaply. The other new separating-weight formulas are different polynomials with proved equivalent zero existence. The fourteen-variable degree is
$18\,930\,435\,180\,664\,062\,500\,000\,002\,752$,
so it is not a low-degree competitor. The nine-positive degree is
$2\,149\,585\,634\,469\,985\,961\,914\,062\,796$.

\section{One coherent coefficient and theorem interface}
To avoid any new compiler-rebinding obligation, the main new families use the entire original compact943 tuple retained by package 30:
\[
 B_*=2^{5^{10}},\quad \beta=B_*-1,\quad
 K_*=D_C+B_*D_R,\quad M_C,\quad M_F=M_F^{\rm native}+B_*-1.
\]
The fourth coefficient is already shifted. The four raw, unsigned little-endian integers occupy 5,312,269 bytes and are included with their recipe and hashes. The smaller $2^{5^9}$ tuples of packages 25 and 27 remain valid competing coefficient choices; retaining the larger established tuple here is deliberate, not a new height record.

The shared coefficient reconstructor reproduces all four source hashes and checks 69,230 support/convolution identities. These finite statements authenticate the recipe and bytes, not its unbounded compiler theorem. That theorem, its canonical completeness properties, and the finitary calibration of the literal K215 blank run remain mathematical dependencies. The resulting proof chain is
\[
\begin{gathered}\text{finite ZFC refutation}\ \Longleftrightarrow\ \text{K215 blank halt}\\
\Longleftrightarrow\ \text{source Diophantine solution}\ \Longleftrightarrow\ \text{new solution}.\end{gathered}
\]
Negation gives the requested ZFC-provable equivalence to $\Con(\mathrm{ZFC})$.
It does not prove either side of that equivalence.

For the Jones families we use package 30's normalized quartic in 31 witnesses. It removes the register $G$ from the old 32-variable system using
\[
 G=q(q-F)-C,\qquad
 L_0G=(q-F)(N-q)-CL_0,\qquad N=q(L_0+1).
\]
Every substitution is polynomial, and the converse restores $G>0$ on positive zeros using the transport equation. The resulting 22 quadratic residuals give a normalized quartic by summing their squares after shifting all 31 positive coordinates by one. The dummy input is fixed to one, so the Jones parameter is $\nu=31$, not 9.

The 402 nonzero factorial-weighted coefficient records and their index recipes are included. Writing $L=5^{32}$ and $s_*=45{,}933{,}757$, the explicit index has
\[
 z_*=2^{s_*},\qquad
 y_*=\sum_{|\alpha|\le4}(z_*+c_\alpha)(2z_*)^{L-\sum_{j=0}^{31}\alpha_j5^j}.
\]
There are $\binom{36}{4}=58{,}905$ positions, including zero coefficients. The bound certificate gives $|c_\alpha|+1<z_*$, and hence
$y_*<(2z_*)^{L+1}$. The integer $z_*$ is materialized; $y_*$ and the larger base-scale constant used below are exact recipes, not materialized payloads. New finalizers are ordinary derivations from the Jones ratio interface, not literal instances of Lean's theorem about its particular \texttt{nineValue} polynomial.

\section{Fourteen unrestricted integer unknowns}
\label{sec:signed14}

We combine package~30's quartic in $31$ positive unknowns, package~27's
fast Jones evaluation, and package~29's signed square-stack and intersective
finalizer. Two new restricted-completeness arguments and a Pell orbit remove
three of package~29's guards. The resulting fixed polynomial has
\[
 \boxed{14\text{ integer unknowns},\quad
 \deg F_{14}=813056\cdot5^{32}+2752,\quad1524\text{ binary operations}.}
\]
A lower-degree alternative has $16$ integer unknowns, degree
$406528\cdot5^{32}+10008$, and $1537$ operations. These are ordinary,
unformalized reductions, not optimality claims or a proof of consistency.
The complete derivation and exact source boundaries are in
\texttt{signed/SIGNED\_PROOF.md}.

\subsection{Fixed index and the exact polynomial}
Use the \emph{whole} compact943 coefficient tuple of package~30 and its
normalized quartic $P_{31}(x,\mathbf n)=\sum_jr_j(\mathbf n+1)^2$.
The dummy $x$ is fixed to $1$. Put $L=5^{32}$ and
\[
 c_\alpha=(4-|\alpha|)!\prod_j\alpha_j!\,[\mathbf X^\alpha]P_{31},
 \qquad z=2^{45933757},
\]
\[
 e(T)=\sum_{|\alpha|\le4}(z+c_\alpha)
 T^{L-\sum_{j=0}^{31}\alpha_j5^j},\qquad y=e(2z).
\]
The $402$ nonzero coefficient records are in
\texttt{signed/upstream/}\allowbreak\texttt{COMPACT\_CONSTANTS.json}. Unlisted coefficients are zero; the index
certificate proves $|c_\alpha|+1<z$.
Consequently $y<(2z)^{L+1}$. Define the additional fixed integer
\[
 \kappa=45933758(L+1)+1,\qquad a=2^\kappa.
\]
Thus $\kappa$ is odd and $a>y$. This enormous coefficient is given by a
bounded recipe, \emph{not materialized}. The same applies to $y$ and the
weighted coefficients. Fixed leaves are not existential variables.

The complete list of unknowns is
\[
 (\varepsilon,g,h,s,w,\phi,i,j,z_1,z_2,z_3,z_4,v,u)\in\mathbb Z^{14}.
\]
Set $b=a\varepsilon^2$, $\omega=w^2$ and define the following
\emph{polynomial abbreviations}, none of which is an extra unknown:
\[
 B=b^5,\quad q=B^L,\quad \ell=\sum_{r=1}^{31}B^{5^r},\quad
 \lambda=\sum_{r=0}^{4L-1}B^r,\quad \theta=B-2z,
 \quad N=q^{16},\quad G=32bB^{5^{31}}.
\]
\begin{align*}
 R={}&\left[g+e(B)q^3+\ell q^5+
 \left(2(e(B)-z\lambda)(1+B+g)^4+\lambda B+\lambda Bq^4\right)q^7\right](N^2-N)\\
 &+\left[q^3-b\ell+\ell+\theta\lambda q^3+(B-2)q^8\right](N^2-1).
\end{align*}
\begin{gather*}
 U=N^2\omega,\quad Y=N^2s,\quad M=RY,\quad P=2M^2U,
 \quad K=R+1+h(P-1),\quad A=M(U+1),\\
 B_1=2R+1,\quad W_1=b\omega,\quad C=B_1+W_1+\phi,
 \quad J=4A-5,\quad Q=3W_1C-2(W_1^2-1),\\
 D=(A^2-1)C^2+1,\quad E=iJDC^2,\quad F=(A^2-1)E^2+1,\\
 G_1=A+F(CD+1-A),\quad H=B_1+jC,\quad I=(G_1^2-1)H^2+1,\\
 a_0=(P^2-1)K^2+1,\quad a_1=DFI,\quad
 m=K^2G-4(C-KY)^2G-8gK^2,\\
 D_0=JF,\qquad T_0=QF+J(H-C).
\end{gather*}
Put $(d_1,d_2,d_3,d_4)=(3g-1,3s-1,3\phi+2,3m-1)$ and
\[
 C_0=1,\qquad B_j'=1+4d_jC_{j-1}^2z_j^2,\qquad
 C_j=C_{j-1}B_j'\ (1\le j\le4),\qquad C_*=C_4.
\]
Finally set
\begin{gather*}
 Z=\varepsilon wi z_1z_2z_3z_4,\quad S=ZD_0^5,\quad r=T_0+D_0v,\\
 T=(r^3-D_0^2r-D_0^3)(r^2+23D_0^2),\quad
 w_1=C_*+1,\quad w_2=w_1(P^2+K^2+2),\quad X=Su-T,\\
 A_0=S^2C_*,\quad A_1=S^2w_1^2a_0,\quad A_2=S^2w_2^2a_1,\\
 V=X^2+A_2-A_0-A_1,\qquad W=V^2+4X^2A_2-4A_0A_1.
\end{gather*}
The requested equation is the single integer polynomial equation
\[
 \boxed{F_{14}=W^2-16V^2X^2A_2=0.}
\]
Its division-free expanded arithmetic DAG is
\texttt{signed/generated/}\allowbreak\texttt{signed14.json}. The exported names
\texttt{residue} and \texttt{normroot} mean $v,u$.

\subsection{The signed Jones interface and a removable guard}
The imported Jones ratio theorem, on its original domain
$b=y+1+\varepsilon_0$, $\varepsilon_0,\phi\ge0$ and
$g,h,s,\omega,i,j\ge1$, identifies quartic solvability with
\[
 a_0,a_1\in\operatorname{Sq}(\mathbb Z),\qquad D_0\mid T_0,
 \qquad m>0.\tag{JR}
\]
For every integer assignment $D_0\ne0$: $J=4A-5\ne0$, and $F=0$
would force $A=0,E=\pm1$, although $E$ would be divisible by $5$.
Replacing a nonzero $i$ by $|i|$ leaves (JR) unchanged since only $E^2$
occurs in its consumers.

Package~29's proof of automatic $h$ positivity does not use $j$.
For completeness, assume $b>y$, $g,s,\omega\ge1$, $\phi\ge0$,
$a_0$ square and $m>0$. The margin gives
\[
 K\ne0,\quad g<4bB^{5^{31}},\quad 4(C-KY)^2<K^2.
\]
Writing $t=5^{31}$, one has $1+B+g<B^{t+2}\le q$ and
$b\ell<q$. In the first bracket for $R$, the potentially negative
term is absorbed by $\lambda Bq^4>2z\lambda(1+B+g)^4$.
Both brackets are positive, so $R\ge2$, $C>0$, and
$P\ge2R^2>R+2$. The margin implies $K>0$. If $h\le0$, then
$1\le K\le R+1<P$. The case $K=1$ would imply $P-1\mid R$;
for $2\le K<P$ the radicand lies strictly between $(PK-1)^2$ and
$(PK)^2$. Both contradict the hypotheses. Hence $h\ge1$.

\paragraph{New Pell-orbit lemma.}
Under those hypotheses and (JR), $A\ge4$, $D\ge A^2$, $F>0$ and
$G_1>A>1$. Since $F\equiv1\pmod D$ and $F\equiv1\pmod J$,
the divisibility gives $H\equiv C\pmod F$. Thus
\[
 G_1\equiv1\pmod D,\quad I\equiv1\pmod D,
 \qquad G_1\equiv A\pmod F,\quad I\equiv D\pmod F.
\]
Therefore $D,F,I$ are positive and pairwise coprime. As $DFI$ is a
square, all three are squares. Put $d=G_1^2-1$, a positive nonsquare,
and $x=\sqrt I>0$. Even if $H<0$, $x+H\sqrt d>0$.
The unit $\eta=G_1+\sqrt d>1$ has norm one, so some positive power
is the identity modulo $CF$ in the integer quadratic algebra.
Multiplying $x+H\sqrt d$ by arbitrarily large such powers yields
\[
 x_n^2-dH_n^2=1,\qquad H_n\equiv H\pmod{CF},\qquad H_n\longrightarrow+\infty.
\]
Choose $H_n>B_1$ and put $j_n=(H_n-B_1)/C>0$. This preserves
$a_1$ being square and $D_0\mid T_0$, and leaves $a_0,m$ unchanged.
Thus the positivity guard on $j$ can be deleted entirely.

\subsection{Cofinal bases and square witnesses: two more deleted guards}
The proof of \texttt{USys\_of\_mem} in pinned
\texttt{Master3.lean}, lines~303--510, chooses an exponent
$k=xy+\sum n_j+1$ only to obtain $k\ge1$, $xy<2^k$, and
$n_j<2^k$ for a fixed quartic witness. The rest of the construction
uses those inequalities, not that equality. Hence it works for every
sufficiently large exponent $k$, including any prescribed parity.
\texttt{NineReduce.reduced\_of\_mem} and
\texttt{NineRatio.ratio\_complete} transport that same base $b=2^k$.
At $x=1$ we may therefore choose $k=\kappa+2r$ large enough, and set
$\varepsilon=2^r$, giving $b=a\varepsilon^2$. Conversely any
nonzero integer $\varepsilon$ gives $b\ge a>y$, so the old ceiling
coordinate $\varepsilon_0=b-y-1$ is nonnegative. This is a cofinal
restriction of the base, not an unjustified square substitution for
$\varepsilon_0$.

There is a second concrete freedom. In pinned \texttt{Ratio.lean},
\texttt{exists\_B25core}, lines~434--442, the canonical witness is
\[
 \omega=2^{2R+1-k},\qquad U=N^2\omega.
\]
The source first proves $k<2R+1$ from $k<2^k=b\le R$, so natural
subtraction does not truncate. Since our $k$ is odd,
$2R+1-k$ is positive and even: this $\omega$ is a positive square.
\texttt{NineRatio.ratio\_complete}, lines~270--279, retains that
same $\omega$ while adding $i,j$ and changing $\phi$. Thus
$\omega=w^2$, $w\ne0$, preserves completeness; soundness follows
from $w^2\ge1$. There is no analogous asserted square property for $s$.

Consequently only $g,s\ge1$, $\phi\ge0$ and $m>0$ require inequality
guards. Nonvanishing of $\varepsilon,w,i$ is sufficient for the other
restrictions. These source-inspected strengthenings are ordinary proofs,
not newly compiled formal theorems.

\subsection{Why the finalizer has exactly the intended zeros}
Every $B_j'$ is $1$ modulo $4$ and coprime to $C_{j-1}$.
If all $z_j\ne0$ and $C_*$ is a square, unique factorization forces
$|B_j'|$ to be squares; their residue modulo $4$ excludes negative
squares. Thus $d_j\ge0$, and $d_j\equiv2\pmod3$ excludes zero.
These are exactly the four required inequalities. Conversely each
positive $d_j\equiv2\pmod3$ is nonsquare. A Pell unit power gives
$u_j^2-d_jv_j^2=1$ with $v_j$ divisible by $2|C_{j-1}|$.
Taking $z_j=v_j/(2C_{j-1})\ne0$ makes every $B_j'$ a square.

The polynomial $\Pi(t)=(t^3-t-1)(t^2+23)$ has no rational root but
has a root modulo every nonzero integer. To see the latter, for an
odd prime other than $23$, either the cubic has a simple root, or
its irreducibility makes Frobenius a three-cycle and its discriminant
$-23$ a nonzero square. Thus one factor has a simple root and lifts.
At $23$ use the cubic root $3$; at $2$ use $-23\equiv1\pmod8$ and
the standard square-root lifting. Apply the Chinese remainder theorem.
For $D\ne0$ the homogeneous expression
\[
 \mathcal R(D,T,v)=((T+Dv)^3-D^2(T+Dv)-D^3)
 ((T+Dv)^2+23D^2)
\]
is nonzero and satisfies
\[
 (D\mid T\ \&\ Z\ne0)\quad\Longleftrightarrow\quad
 \exists v:\ ZD^5\mid\mathcal R(D,T,v).
\]
The forward direction applies intersectivity after $T=Dq$.
For the converse, if $v_p(D)=\alpha>\beta=v_p(T)$, the two factors
have valuations $3\beta,2\beta$, contradicting divisibility by $D^5$.
The case $Z=0$ would contradict nonvanishing. This supplies exactly
$D_0\mid T_0$ and all nonzero conditions via $S\mid T$.

The new linear weight is valid even before recovering signs:
$C_*\equiv1\pmod4$ implies
\[
 |C_*+1|>\sqrt{|C_*|},\qquad
 \sqrt{|a_0|}\le P^2+K^2+1.
\]
For negative $C_*=-t$, necessarily $t\ge3$ and
$(t-1)^2-t=t^2-3t+1>0$. Hence
\[
 |w_2|>\sqrt{|C_*|}+|w_1|\sqrt{|a_0|}.
\]
For complex roots $r_0^2=C_*,r_1^2=a_0,r_2^2=a_1$, multiplication gives
\[
 F_{14}=\prod_{\sigma\in\{-1,1\}^3}
 \bigl[X-S(\sigma_0r_0+\sigma_1w_1r_1+\sigma_2w_2r_2)\bigr].
\]
If $S=0$, this is $T^8\ne0$. Otherwise a zero makes one encoded
root sum rational. The strict absolute-weight bounds separate all
sign choices on nonzero roots, including imaginary roots. Every
Galois automorphism of their multiquadratic field must therefore fix
each root, so all three roots are rational, hence integers. It follows
that $S\mid T$ and $C_*,a_0,a_1$ are integer squares. Conversely these
conditions yield a zero by taking
$u=T/S+r_0+w_1r_1+w_2r_2$. Zero roots and repeated square classes
cause no exception.

Combining this exact elimination with the square stack, homogeneous
merger, recovered $h$, sign change of $i$, and Pell orbit for $j$
recovers the original Jones witness. Conversely the cofinal-base and
square-$\omega$ constructions, then the stack and finalizer, produce
an integer zero from any quartic witness. The imported compiler and
literal K215 calibration give
\[
 \mathrm{ZFC}\vdash
 \bigl(\nexists\mathbf t\in\mathbb Z^{14}:F_{14}(\mathbf t)=0\bigr)
 \ \longleftrightarrow\ \operatorname{Con}(\mathrm{ZFC}).
\]

\subsection{Lower-degree alternative and exact complexity}
For $F_{16}$ keep $b=y+1+\varepsilon$, $\omega=w$ and use six
stack discriminants
\[
 3\varepsilon+2,\ 3g-1,\ 3s-1,\ 3w-1,\ 3\phi+2,\ 3m-1,
 \qquad Z=i\prod_{r=1}^6z_r.
\]
The same finalizer and signed linear weights apply. Its proof uses
only the new $j$-orbit and weight improvement, not either cofinal
specialization. The intermediate $15$-variable construction restricts
only $b$; its $1530$ operations and degree $813056L+4184$ are dominated
by $F_{14}$.

The fast evaluator shares fifth powers, evaluates the background digit
sum by the complete homogeneous polynomial $h_4$, and uses fixed
multivariate Horner evaluation for sparse corrections. In particular,
$1+t+t^2+t^3+t^4=1+(t^2+t)(t^2+1)$ saves a gate when $t^2$ is shared.
All scalar multiplications and squares are charged; fixed coefficient
construction is excluded, as in the comparison convention.
For $F_{14}$,
\[
\begin{array}{c|c@{\qquad}c|c}
 R&470L&P&1900L+4\\
 K&1900L+5&a_1&49280L+98\\
 m&4442L+14&C_*&4442L+133\\
 S&68750L+152&w_2&8242L+143\\
 A_2&203264L+688&F_{14}&813056L+2752
\end{array}
\]
Here $\deg A_2>\deg X^2,\deg A_0,\deg A_1$, so the successive
leading terms of $V,W,F_{14}$ are $A_2,A_2^2,A_2^4$.
A nonzero symbolic diagonal leading coefficient, certified modulo
$1000003$ as a scalar times a power of the positive fixed $a$, proves
these are exact degrees. For $F_{16}$ one obtains
$\deg A_2=101632L+2502$.

The delivered generator and checker reproduce the DAGs, counts, and
degrees. Diagnostics include $216$ actual-coefficient modular circuit
comparisons, $1000$ eight-factor norm evaluations, $840$ coprimality
cases, $120$ negative-to-positive Pell orbits, and intersectivity tests
for every modulus $1$ through $1200$. These checks supplement, and do
not replace, the universal proofs above.

\section{Positive norms: degree 1061 and a 109-operation evaluator}
\label{sec:positive-wave5}

All unknowns in this section are strictly positive.  There are exactly twelve:
$J,F,\alpha,z,h,i,T,s,w,\eta,\zeta,t$.  Every other name is a polynomial
abbreviation.  All constructions use the complete source-30 Wang943 tuple
$(\beta,K_0,M_C,M_F)$, with $\beta=2^{5^{10}}-1$ and the shifted high mask
$M_F$.  The shared coefficient rebuild reproduces all four selected bytestrings.

\paragraph{Common interface.}
Write $\omega=w$, except in the lowest-degree family, where $\omega=w+q$.
Define
\begin{gather*}
 q=\beta J+1,\quad C_w=q-F-\alpha,\quad X=\omega q,\quad Y=sq^3,\quad E=XY,\\
 k=\eta+\zeta,\quad c=kY+\eta,\quad a=E+Y,\quad H=4a+3,\quad\Delta=a^2+H,\\
 \ell=EkY,\quad A=1+\ell(\ell+k),\quad G=1+\Delta c^2,\quad b_0=X+ac,\\
 R=((q-1)(q-F)+\alpha)(q^2-1)+(M_C+qM_F)J,\\
 N_k=k-hE-R,\qquad N_t=(K_0+\omega)C_w+q-F-z(q-1),\\
 I=ic^2,\quad D=1+\Delta I^2,\quad Q=(\Delta I)^2,\quad d=Q-1,\\
 v=cT,\quad b=c+RD,\quad e=Q(Dv^2+b^2)-1,\quad B_1=2Qvb.
\end{gather*}
Use one of
\[
 C_1=E+3(X+Y),\quad C_2=E+2(X+Y)+sq^2,\quad
 C_3=E+2(X+Y)+sq,
\]
where $C_3$ is used only with $\omega=w+q$.  Put
\[
 M=I(2T+RC_jic),\qquad L=1+HM,\qquad W=Lb_0+Ht.
\]
The new norm is the following integer polynomial, with no additional witnesses:
\[
 \mathcal N=\prod_{g=\pm\sqrt G,\,\tau=\pm\sqrt A,\,f=\pm\sqrt D}
 \left[d(W-Lg-HM\tau)^2-H^2e+H^2B_1f\right].
\]
The presentations are $N_k=N_t=1,\ \mathcal N=0$, and the single equation
\[
 P=N_kN_t(1+\mathcal N)-1=0.
\]

\paragraph{A half-sized carrier and the new restriction.}
Source 24 used $L=1+2HM$.  Its uniform conjugate estimate
$|y|<M\sqrt A$ remains valid here, where
$dy^2=e-B_1f$.  Hence all roots of the inner polynomial are
$x=M\tau+y$ with $|x|<2M\sqrt A$.
Before any root recovery, the authentic weak interface gives
$q\ge16$, $X\ge16$, $Y\ge4096$, $E>R+2$, $k>E$, and $c>kY$.
Consequently $G>A$ and
$b_0=\ell+E\eta+Yc+X>\ell+k>\sqrt A$.
For the proposed smaller carrier, the negative main-root branch would give
\[
 x=\frac{W+L\sqrt G}{H}
   =t+(M+H^{-1})(b_0+\sqrt G)>2M\sqrt A,
\]
a contradiction.  If $\sqrt G$ were irrational, its negative conjugate
would also be an inner root; thus the main root is a positive integer.
Conjugation next forces $\tau$ rational, then $y$ and $f$ rational.
The inherited rational Pell saturation lemma makes $y$ integral.
Reduction modulo $H$, using $L\equiv1\pmod H$, restores
$g=b_0+H\gamma$, with $\gamma>0$.  The unchanged signed-quotient compiler
interface then supplies the accepting tableau.  Conversely a canonical
positive witness extends by $t=y+M\tau+L\gamma>0$.

For the degree-1061 version, $Y/X=sq^2/(w+q)<sq$, so $C_3$ satisfies the
same source-24 bound that previously used $sq^2$.  Completeness is not
assumed: after authentic root recovery, $X=\omega q=2^R$ and $R>3q$,
so $\omega>q$.  Thus $w=\omega-q$ is strictly positive.  Soundness of
$\omega=w+q$ is immediate.  This coordinate change is not literal identity
of polynomials at the same supplied $w$.

\paragraph{Cheaper trace--remainder evaluation.}
Source 30's evaluator applies to every carrier above.  Let
\begin{align*}
 n&=(HM)^2A,&d_n&=dn,&H_Q&=H^2Q,\\
 H_P&=H_QDv^2,&H_U&=H_Qb^2,&H_e&=H_P+H_U-H^2,&j&=4H_PH_U,\\
 T_r&=2d(W^2+L^2G),&K_r&=d^2(W^2-L^2G)^2,\\
 a_t&=T_r-2(d_n+H_e),&b_t&=(d_n-H_e)^2-j-K_r,\\
 \bar v&=16d_nj-a_t(a_tT_r+2b_t),&u_r&=b_t^2-a_t^2K_r.
\end{align*}
Then
\[
 \mathcal N=K_r\bar v^2-u_r(T_r\bar v-u_r).
\]
Sharing $H_Q$ saves one multiplication; the final Horner expression saves
another.  Comparing its two terms directly saves one further operation
for a system.  These are exact identities, including degenerate quadratic
specializations.  Also $\mathcal N$ is a square modulo $4$, so the product
equation forces $\mathcal N=0$ and $N_k=N_t\in\{1,-1\}$; the inherited
weak-sign compiler restores the positive unit sign.

\begin{center}
\begin{tabular}{lrrrr}
\hline
Choice & Product degree & Product ops & System degree & System ops\\
\hline
Source 30, unchanged polynomial &1081&109&1072&104 / 103\\
$C_1$, half carrier &1069&117&1060&112 / 111\\
$C_2$, half carrier &1065&119&1056&114 / 113\\
$C_3$, shifted $w$ &1061&120&1052&115 / 114\\
\hline
\end{tabular}
\end{center}
The second system count uses the proved shift $z_{\rm old}=z+1$, replacing
$N_t=1$ by $(K_0+\omega)C_w=F+z(q-1)$.  All binary $+,-,\times$ gates,
including multiplication by numerals, are charged; comparisons are free.

\paragraph{Exact degree and proof boundary.}
For the $C_2$ norm, the leading homogeneous form is
\[
 2^{46}i^{48}(\eta+\zeta)^{96}w^{74}s^{176}
 (\beta J)^{646}(\beta J-F)^{16},
\]
of degree 1056.  Four of the sixteen linear conjugate factors have degree
35, two have degree 36, and ten have degree 38; the multiplier is $d^8H^{16}$.
For $C_3$, the first four degrees drop to 34.  Their complete leader contains
\[
 X_0^2Y_0^2(\eta+2\zeta)+(\eta+\zeta)s^3q_0^8w,
\]
which is nonzero; retaining the tied $-E_0^2\eta$ term is essential.
This proves norm degree 1052, and $N_kN_t$ adds nine degrees.
The standalone proof supplies all separation estimates and leading forms.
Independent exact conjugate products and complete finite-field ray checks
corroborate the circuits, not the unbounded compiler theorem.
These are ordinary reductions formalizable in ZFC, not a consistency proof,
a global optimality claim, or a new kernel-checked endpoint.

\section{A sharper equation in nine natural or positive unknowns}
\label{sec:nine-completed-square}

The new nine-variable polynomial has exact total degree
\[
 \boxed{92324\cdot5^{32}+296
       =2149585634469985961914062796.}
\]
Its supplied integer arithmetic circuit uses \textbf{1505 binary operations}
 over $\mathbb N=\{0,1,\ldots\}$, or \textbf{1500 operations} after the
coordinate translation to strictly positive integers.  The earlier source27 positive circuit still uses fewer operations (1480), but at much higher degree; the new form does not dominate it on cost alone. Source~30 used
$120724\cdot5^{32}+366$ and 1840 operations.  The new construction retains
source~30's complete compact943 coefficient tuple and 31-witness quartic,
combines them with source~28's signed-divisor finalizer, and refines its
majorants.  This is an ordinary mathematical synthesis of the stated
imported theorems, not a newly kernel-checked endpoint or a minimality claim.

\paragraph{The fixed instance.}
The four ports are the unchanged compact943 tuple
$(\beta,K_0,M_C,M_F)$, with native radix $2^{5^{10}}$ and
$M_F=M_{F,\mathrm{native}}+\beta$.  Let $r_1,\ldots,r_{22}$ be the
literal quadratic residuals in \texttt{nine/source/30-quartic31/arithmetic/}.
Their sum of squares after shifting all 31 positive coordinates by one is
$P_{31}(x,\mathbf n)$, independent of the dummy input $x$.
It has degree four and $P_{31}(x,\mathbf0)\ge1$ because residual~5 equals
one at the all-ones tuple.  Source~30's exact elimination of one register
transfers the inherited compiler theorem to this parent.

Put $\nu=31$, $L=5^{32}$ and
$c_\alpha=(4-|\alpha|)!\prod_j\alpha_j!\,[\mathbf x^\alpha]P_{31}$.
All 402 nonzero coefficient recipes are supplied.  The recorded bit lengths
of the four authentic ports give, for every recipe
$c_\alpha=\sum_t a_t\boldsymbol\kappa^{\mathbf e_t}$,
\[
 |c_\alpha|+1\le
 1+\sum_t|a_t|2^{\sum_j e_{tj}\operatorname{bits}(\kappa_j)}
 <2^{45933757}.
\]
Thus $z=2^{45933757}$ and
$y=\sum_{|\alpha|\le4}(z+c_\alpha)(2z)^{L-\sum_j\alpha_j5^j}$
give the explicit admissible Jones index.  The baseline includes all
58,905 exponent positions.  These are bounded coefficient recipes, not
additional unknowns or an unspecified enumeration index.

\paragraph{The new majorants.}
Use the unchanged Jones packing and ratio expressions, fully listed in
\texttt{nine/PROOF.md}, on the eight natural outer coordinates.  In
particular
\[
 A=R\sigma,\quad \sigma=Y(U+1),\quad r=bw+\phi,\quad
 C=2R+r+1,\quad H=(j+1)C-r,
\]
\[
 D=(A^2-1)C^2+1,\quad E=i(4A-5)DC^2,\quad
 F=(A^2-1)E^2+1,\quad G=FDC+F-A(F-1).
\]
Source~28's unconditional congruence
$R\equiv(b-1)\ell\pmod q$, with $0<(b-1)\ell<q$, proves $R\ne0$.
Consequently $|A|\ge2$, $P\ge2$, $K\ge1$, $AE\ge0$ and $D,F\ge1$;
no margin or square condition is assumed.  Replace the earlier majorants by
\begin{align*}
 M_D&=AC+\sigma(r+1)^2+1,& M_F&=AE+1,\\
 T_{CH}&=CH+r^2,&
 M_I&=F(DT_{CH}+2H^2+A^2)-A^2+1,\\
 V_0&=PK+1,&V_1&=M_DM_FM_I+1.
\end{align*}
Indeed $M_D-AC>0$ and
\[
 M_D+AC=\sigma\{4R^2+2R(r+1)+(r+1)^2\}+1>0,
\]
so $M_D\ge|AC|+1\ge\sqrt D$.  Also $M_F\ge\sqrt F$, and
\[
 T_{CH}+CH=C^2+(C-r)^2+2jC^2\ge0,
 \qquad T_{CH}-CH=r^2\ge0.
\]
Thus $T_{CH}\ge|CH|$.  The decomposition of $G$, together with
$|H|\le H^2$ and $|AH|\le A^2+H^2$, gives
$M_I\ge|GH|+|H|+1\ge\sqrt{|(G^2-1)H^2+1|}$.
For $a_0=(P^2-1)K^2+1$ and $a_1=DF((G^2-1)H^2+1)$,
we obtain $V_t\ge2$ and $\sqrt{|a_t|}\le V_t-1$.
The proof includes negative $R$, $C=0$ and $H=0$.

\paragraph{The signed finalizer.}
Let $u=JF\ne0$, $v=QF+J(H-C)$, with $J=4A-5$, and let $m$ be the
unchanged Jones margin.  Put $M=Q^2+(H-C)^2+1$; then
$|v/u|\le|Q|+|H-C|\le M-1$.  With the ninth unknown $n\in\mathbb N$,
\[
 \Omega=u(2m-1),\qquad
 \Theta=u(n+M)+v-\Omega(2M+V_0V_1),
\]
\[
 \mathcal F_9=
 \bigl(\Theta^2-\Omega^2(a_0+V_0^2a_1)\bigr)^2
       -4\Omega^4V_0^2a_0a_1.
\]
Source~28's four-factor radical argument proves that
$\exists n\in\mathbb N:\mathcal F_9=0$ is equivalent to
$a_0,a_1$ being squares, $u\mid v$ and $m>0$, for either sign of $u$.
For clarity, a vanishing factor yields integer roots $r_t^2=a_t$:
if $s=r_0+V_0r_1\in\mathbb Q$ is nonzero, use
$2sr_0=s^2+a_0-V_0^2a_1$; if $s=0$, the bounds force $a_1=a_0=0$.
Then $Z=2M+V_0V_1+r_0+V_0r_1\ge2M+1$, and
$n+M+v/u=(2m-1)Z$ gives soundness.  Conversely choose
$n=(2m-1)Z-M-v/u\ge2$.  The unchanged Jones reduction, parent compiler
and literal K215 calibration therefore give, provably in ZFC,
\[
 \bigl(\nexists\mathbf x\in\mathbb N^9:
          \mathcal F_9(\mathbf x)=0\bigr)
 \quad\Longleftrightarrow\quad\operatorname{Con}(\mathrm{ZFC}).
\]

\paragraph{Why the degree drops.}
The leading terms of $V_1^2$ and $a_1$ now cancel.  This is resolved by an
exact telescoping certificate, not a naive gate-degree estimate.  With
$t=\sigma(r+1)^2$,
\[
 M_D^2-D=2ACt+t^2+2AC+2t+C^2
\]
has degree $1110L+8$ and diagonal leader $4$.  In
\begin{align*}
 V_1^2-a_1={}&(M_D^2-D)M_F^2M_I^2
 +D(M_F^2-F)M_I^2\\
 &+DF(M_I^2-I)+2M_DM_FM_I+1,
\end{align*}
the first term uniquely dominates: its degree is $24170L+72$ and leader
$2^{46}$.  Writing $E_0=u(n+M)+v-2\Omega M$, the inner norm is exactly
\[
 E_0^2-2E_0\Omega V_0V_1+\Omega^2V_0^2(V_1^2-a_1)-\Omega^2a_0.
\]
The third term uniquely dominates with degree $46162L+148$ and leader
$2^{98}$.  The outer square therefore has degree $92324L+296$ and
leader $2^{196}>0$; the final subtracted term has smaller degree.

\paragraph{Executable representation and checks.}
The shared radix-five ladder and complete-homogeneous $h_4$ identity from
source~27 evaluate the 58,905-position digit baseline without expanding it.
Exact coefficient aliases leave 78 nonliteral coefficient leaves besides
$z,y$.  The identity
$1+t+t^2+t^3+t^4=1+(t^2+t)(t^2+1)$, deterministic multivariate Horner
ordering and further distributive rewrites give the stated operation count.
The supplied standard-library generator and verifier rebuild all 402
coefficient recipes, prove the complete sparse digit identity, check the
actual coefficient hashes and index bound, audit the cancellation-free
degree circuit and compare actual-instance evaluations modulo four primes.
Finite tests supplement the universal proof; they do not decide
solvability or recheck the complete semantic chain in a proof-assistant
kernel.

\section{Integer-domain synthesis}
Two combinations improve different parts of the integer frontier. All their
coefficient ports select the same complete compact Wang943 tuple, with
$B=2^{5^{10}}$ and the shifted high-mask port. Historical source bindings are
not spliced into these new instances.

\paragraph{A quartic in fifty unrestricted integers.}
Source~30 removes the positive packing register $G$ by the identity
\[
 L\bigl(q(q-F)-C\bigr)-\bigl((q-F)(N-q)-CL\bigr)
   =(q-F)\bigl(q(L+1)-N\bigr).
\]
This identity is equally valid in source~27's signed odd-scale lift. Deleting
$G$ there, while retaining $i=2r+1$, $s=2v+1$, $w=u^2$ and six three-variable
positive guards, gives fifty unrestricted integers and 29 quadratic equations.
Restoring $G=q(q-F)-C$ is unique and precedes every sign-recovery argument;
thus no positive-$s$ assumption is imported prematurely. The sum of squares
has exact degree four, since its $\tau^4$ coefficient is one. Its comparison
system has 111 binary operations, and the explicit quartic evaluator has 197.
The expanded polynomial has 414 symbolic terms, including coefficient ports.

\paragraph{A cheaper degree-1936 polynomial in twenty-four integers.}
Keep the signed ratio-norm notation $I=ic^2$, $v=cT$, $b=c+RD$,
$G=1+\Delta c^2$. The improved weight is
\[
 M=2(5Rac\,iI+v^2),\qquad L=1+HM.
\]
The inherited sign-safe bound strengthens to
\[
 \frac{|y'|}{\sqrt G}
 <iI+v^2+2+4Rac\,iI
 <5Rac\,iI+v^2=M/2,
\]
because $(Rac-1)iI>2$ already holds before root recovery. Both signs of
$s,i,H$, and the $w=0$, $T=0$ boundaries, remain covered. With
$S=X+ac+Ht$, take the eight-factor norm of
\[
 d(S-Lu-2HM\tau)^2-H^2e+H^2B_{\rm aux}f,
 \qquad u^2=G,\quad\tau^2=A,\quad f^2=D.
\]
Source~30's trace/product remainder method evaluates this norm. Sharing
$HQ=H^2Q$ and using a Horner final norm further reduces its circuit.
The product $N_kN_t(1+N)-1$ now costs 161 operations at exact degree1936.
The three-equation presentation costs 156 at maximum degree1920, using a
direct final comparison; its SOS costs 164 at degree3840.
The uniform leading certificate is
\[
 N_*=(99\cdot899)^4d_*^8H_*^{16}b_*^{16},
 \qquad \deg N=1920.
\]
This replaces the previous degree1936 evaluations at 167/168 operations.

\paragraph{Verification and scope.}
The generators deliver every binary gate and every quadratic residual.
Independent checks establish the packing syzygy, the quartic expansion,
208 rank-eight conjugate-product comparisons, and exact finite-field degree
certificates using the actual coefficient tuple. Full sign proofs and source
provenance are in \texttt{lowdegree/PROOF.md}. The compiler and calibrated
K215 endpoint remain inherited mathematical dependencies. These constructions
assert neither consistency nor global minimality, and are not new end-to-end
proof-assistant certificates.


\section{Validation, attribution, and limitations}
Every new family includes a generator, explicit arithmetic circuits or sparse residuals, a mathematical proof, and replay checks. The final verification receipt records the actual completed stages and byte-identical regenerated artifacts. Separate independent reviews address the domain-restoration lemmas and root separation. Corruption controls are designed to reject a changed coefficient payload, a changed circuit, and a missing source dependency. They do not test or decide whether a genuine ZFC-refutation witness exists.

The full source comparisons and proof notes identify inherited facts separately from new arguments. The symbolic identities establish polynomial equality where that is claimed; modular evaluations are diagnostics, not proofs of equality. Exact-degree claims include nonzero leading-term certificates rather than merely structural upper bounds. All original source snapshots retain their provenance and original tuple metadata; only newly generated packets select the shared binding.

This archive does not replace an endpoint build. The published Lean/Rocq K215 endpoint and the inherited native-window/Pell compiler are used through their identified ordinary mathematical calibration. New cofinal-base and finalizer arguments have not been compiled into either proof assistant. Package 27's earlier licence omission is not repeated: the Wade licence and K215 notices are included alongside the repository licence and source attribution. Source scripts or snapshots omitted during repository intake are not misidentified as defects of the original sealed archives.

P67 remains the lowest-operation single equation in the reviewed collection. The older low-degree and low-operation systems remain useful when they are not dominated. Fourteen unrestricted integers is a concrete improvement over the reviewed seventeen-integer construction, not a global lower bound or a claim to supersede abstract representation theorems. All claims concern the stated domains and the coherent fixed coefficient tuple.

\end{document}
